% TOP_PROBLEM: {"id": 2307049, "title": "Research Problems in Function Theory — Problem 7.49", "category_id": 8, "category": {"id": 8, "name": "analysis", "display_name": "Analysis", "description": "Limits, continuity, calculus, and function theory.", "slug": "analysis", "order_index": 8, "created_at": "2026-07-31T15:26:25.671Z"}, "status": "open"}
% FIRST_POSTED: null
% ENTRY_KIND: statement_only
% Imported from: batch05.tex; UnsolvedMath ID 2307049
\documentclass[11pt]{article}
\usepackage{amsmath,amssymb,mathtools}
\usepackage{fontspec,unicode-math}
\setmainfont{Latin Modern Roman}
\setmathfont{Latin Modern Math}
\usepackage{xurl,hyperref}
\newcommand{\sref}[2]{\hyperlink{source:#1:#2}{[S#2]}}
\newcommand{\cataloguescope}[1]{\paragraph{Mathematical statement.}#1}
\begin{document}
% UnsolvedMath ID 2307049; source catalogue code AMR-022-7049
\setcounter{section}{2307048}
\section{Research Problems in Function Theory — Problem 7.49}\label{2307049}
{\small\sffamily UnsolvedMath ID 2307049 · Analysis}
\subsection{Definitions and mathematical statement}

\paragraph{Shared notation.}$\mathbb N=\{1,2,\ldots\}$, and $\mathbb Z,\mathbb Q,\mathbb R,\mathbb C$ denote the integers, rationals, reals and complex numbers. Any other conventions are specified locally.


Fix $n\ge3$, an integer $m\ge2$ and $p\in(1,\infty)\setminus\{2\}$. A real polynomial $u$ is homogeneous of degree $m$ if $u(tx)=t^m u(x)$ for every $t>0$. Write $Du$ for its gradient, $D^2u$ for its Hessian and $\Delta u$ for the sum of its second partial derivatives. The polynomial form of the $p$-Laplace equation is
\[|Du|^2\Delta u+(p-2)\langle D^2u\,Du,Du\rangle=0.\]
At points where $Du\ne0$ this is equivalent to $\operatorname{div}(|Du|^{p-2}Du)=0$; the displayed polynomial identity also makes sense at critical points.

\paragraph{Statement recorded in the supplied source.}
Must every real homogeneous polynomial of degree $m\ge2$ satisfying this identity throughout $\mathbb R^n$ be identically zero, for every $n\ge3$ and every $1<p<\infty$, $p\ne2$? \sref{2307049}{2}
Tkachev's Conjecture1.1 explicitly states Lewis's full question and fixes the intended differential equation, rather than the malformed gradient-dot expression printed in the problem list. The paper records the quadratic case and proves the cubic case without restricting the degree of the question. \sref{2307049}{3}
\subsection{Short English statement}
Can a nonzero real homogeneous polynomial of degree at least two solve the nonquadratic p-Laplace equation in dimension at least three?
\subsection{Sources}
\begin{description}
\item[\hypertarget{source:2307049:1}{S1}] ULAM.AI and the UnsolvedMath Contributors, \emph{UnsolvedMath: A Curated Collection of Open Mathematics Problems}, record 2307049 (AMR-022-7049). CC BY4.0. \url{https://huggingface.co/datasets/ulamai/UnsolvedMath}.
\item[\hypertarget{source:2307049:2}{S2}] W. K. Hayman and E. F. Lingham, \emph{Research Problems in Function Theory} (2018), Chapter7, Problem7.49 and its complete update. \url{https://arxiv.org/abs/1809.07200}.
\item[\hypertarget{source:2307049:3}{S3}] Vladimir G. Tkachev, \emph{On the non-vanishing property for real analytic solutions of the $p$-Laplace equation} (2015), Section1, equation(1.1), Conjecture1.1 and Theorem1.2, pages1--2. \url{https://arxiv.org/abs/1503.03234}.
\end{description}
\label{end:2307049}
\end{document}
